% ============================================================
%  EXAMEN — TEMA 1 (PRIMER LAPSO)
%  Fracciones Equivalentes, Operaciones Combinadas,
%  Potenciación y Ecuaciones en Q
%  Curso: 2do Año — CELAM
%  Autor: Ing. Gabriel Astudillo
%  Basado en guia_t1l1.tex (evaluación individual de 20 puntos)
%  Nota: SOLO el enunciado. No incluye clave de respuestas ni
%  indicadores de dificultad (sin estrellitas). No hay problemas
%  de planteamiento: solo ejercicios de cálculo.
% ============================================================

\documentclass[12pt, letterpaper]{article}

% ── Paquetes ─────────────────────────────────────────────────

\usepackage{mathwizards-guia}
\mwguia{Examen — Tema 1}{Ing. Gabriel Astudillo $\cdot$ 2do Año}

% ── Comandos propios de este examen ──────────────────────────
\newcommand{\Q}{\mathbb{Q}}

\begin{document}
% ============================================================

% ── Portada / Título ─────────────────────────────────────────
\mwportada{Examen — Tema 1}{Números Racionales ($\Q$)}{Fracciones Equivalentes $\cdot$ Operaciones Combinadas $\cdot$ Potenciación $\cdot$ Ecuaciones}

\vspace{4pt}
\begin{cuadroinfo}
  \small
  \textbf{Nombre:} \rule{0.55\linewidth}{0.4pt} \quad \textbf{Fecha:} \rule{0.15\linewidth}{0.4pt} \\[6pt]
  \textbf{Tiempo disponible:} 90 minutos \quad$\bullet$\quad \textbf{Puntaje total:} 20 puntos \\[4pt]
  \textbf{Estructura:} 4 secciones $\cdot$ 12 ejercicios. \\[4pt]
  \textbf{Instrucciones:}
  \begin{enumerate}[leftmargin=*]
    \item Lee cada ejercicio con calma antes de resolverlo.
    \item Muestra \emph{todos} tus pasos: se evalúa el procedimiento, no solo la respuesta final.
    \item No se permite el uso de calculadora.
    \item Respeta la \emph{jerarquía} de las operaciones y los signos de agrupación.
    \item Simplifica siempre hasta la fracción irreductible y no dejes exponentes negativos.
    \item Escribe con letra clara y revisa tu examen antes de entregarlo.
  \end{enumerate}
\end{cuadroinfo}

\vspace{8pt}

% ============================================================
\section{Fracciones Equivalentes (4 puntos)}
% ============================================================

\begin{enumerate}[leftmargin=*]

  \item \textbf{(2 puntos)} Halla el valor de $x$ en cada igualdad.
        \begin{enumerate}[label=\alph*), leftmargin=*]
          \item $\dfrac{3}{7} = \dfrac{x}{56}$
          \item $\dfrac{45}{x} = \dfrac{9}{11}$
          \item $\dfrac{x + 2}{5} = \dfrac{21}{15}$
          \item $\dfrac{4x - 1}{9} = \dfrac{7}{3}$
        \end{enumerate}

  \item \textbf{(1 punto)} Determina si cada par de fracciones es \emph{equivalente}. Justifica simplificando.
        \begin{enumerate}[label=\alph*), leftmargin=*]
          \item $\dfrac{18}{24}$ y $\dfrac{27}{36}$
          \item $\dfrac{14}{35}$ y $\dfrac{22}{55}$
          \item $\dfrac{84}{126}$ y $\dfrac{55}{99}$
        \end{enumerate}

  \item \textbf{(1 punto)} Completa los valores que faltan para que todas las fracciones sean equivalentes.
        \begin{enumerate}[label=\alph*), leftmargin=*]
          \item $\dfrac{2}{5} = \dfrac{\square}{35} = \dfrac{26}{\square} = \dfrac{\square}{100}$
          \item $\dfrac{\square}{72} = \dfrac{10}{\square} = \dfrac{5}{6}$
        \end{enumerate}

\end{enumerate}

% ============================================================
\section{Operaciones Combinadas (6 puntos)}
% ============================================================

\begin{enumerate}[leftmargin=*, resume]

  \item \textbf{(2 puntos)} Calcula respetando la jerarquía de las operaciones.
        \begin{enumerate}[label=\alph*), leftmargin=*]
          \item $\dfrac{3}{4} + \dfrac{2}{3} \cdot \left(\dfrac{5}{6} - \dfrac{1}{2}\right) \div \dfrac{7}{9}$
          \item $\dfrac{5}{8} - \dfrac{1}{4} \div \dfrac{3}{2} + 1$
        \end{enumerate}

  \item \textbf{(2 puntos)} Calcula respetando la jerarquía y los signos de agrupación.
        \begin{enumerate}[label=\alph*), leftmargin=*]
          \item $\left[\left(\dfrac{5}{8} - \dfrac{1}{4}\right) \div \dfrac{3}{4} + \dfrac{2}{3}\right] \cdot \dfrac{9}{10}$
          \item $\left\{1 - \left[\dfrac{2}{5} + \left(\dfrac{1}{2} - \dfrac{1}{3}\right)\right]\right\} \div \dfrac{7}{30}$
        \end{enumerate}

  \item \textbf{(2 puntos)} Calcula, resolviendo primero los numeradores y denominadores.
        \begin{enumerate}[label=\alph*), leftmargin=*]
          \item $\dfrac{\dfrac{2}{3} + \dfrac{1}{4}}{\dfrac{5}{6} - \dfrac{3}{8}} \div \left(1 - \dfrac{2}{5}\right)$
          \item $\left(\dfrac{7}{4} \div \dfrac{7}{8} - \dfrac{3}{2}\right) \cdot \dfrac{6}{5} + \dfrac{1}{10}$
        \end{enumerate}

\end{enumerate}

\newpage

% ============================================================
\section{Potenciación en $\Q$ (5 puntos)}
% ============================================================

\begin{enumerate}[leftmargin=*, resume]

  \item \textbf{(2 puntos)} Expresa como una \emph{sola} potencia. Observa los factores repetidos.
        \begin{enumerate}[label=\alph*), leftmargin=*]
          \item $\left(\dfrac{2}{3}\right)^4 \cdot \left(\dfrac{2}{3}\right)^3 \cdot \left(\dfrac{2}{3}\right)^5$
          \item $\left(-\dfrac{7}{5}\right)^9 \div \left(-\dfrac{7}{5}\right)^6$
          \item $\left[\left(\dfrac{3}{4}\right)^3\right]^4$
          \item $\left(\dfrac{5}{2}\right)^7 \div \left(\dfrac{5}{2}\right)^7$
        \end{enumerate}

  \item \textbf{(1,5 puntos)} Calcula el valor exacto, cuidando el signo del resultado.
        \begin{enumerate}[label=\alph*), leftmargin=*]
          \item $\left(-\dfrac{2}{5}\right)^4$
          \item $\left(-\dfrac{1}{3}\right)^5$
          \item $-\left(\dfrac{3}{2}\right)^2$
          \item $\left(\dfrac{4}{9}\right)^{-2}$
        \end{enumerate}

  \item \textbf{(1,5 puntos)} Simplifica aplicando las propiedades de las potencias y escribe sin exponentes negativos.
        \begin{enumerate}[label=\alph*), leftmargin=*]
          \item $\left(\dfrac{27\,a^5 b^3}{9\,a^2 b^7}\right)^{-2}$
          \item $\left(\dfrac{3 a^4 b^2}{2 c^3}\right)^{-2} \cdot \left(\dfrac{4 c^6}{9 a^8 b^4}\right)$
          \item $\dfrac{\left(\dfrac{2}{3}\right)^5 \cdot \left(\dfrac{3}{4}\right)^3}{\left(\dfrac{1}{2}\right)^4 \cdot \left(\dfrac{4}{9}\right)^2}$
        \end{enumerate}

\end{enumerate}

% ============================================================
\section{Ecuaciones en $\Q$ (5 puntos)}
% ============================================================

\begin{enumerate}[leftmargin=*, resume]

  \item \textbf{(2 puntos)} Resuelve cada ecuación eliminando primero los denominadores. Comprueba tu solución.
        \begin{enumerate}[label=\alph*), leftmargin=*]
          \item $\dfrac{2x - 1}{3} - \dfrac{x + 2}{4} = \dfrac{x - 3}{2}$
          \item $\dfrac{3x + 2}{5} + \dfrac{x - 1}{2} = \dfrac{5x + 3}{10}$
        \end{enumerate}

  \item \textbf{(2 puntos)} Resuelve cada ecuación y verifica tu resultado.
        \begin{enumerate}[label=\alph*), leftmargin=*]
          \item $\dfrac{x}{2} - \dfrac{x - 1}{3} = x - \dfrac{2x + 3}{6}$
          \item $\dfrac{x + 3}{4} - \dfrac{2x - 1}{6} = \dfrac{x - 2}{3}$
        \end{enumerate}

  \item \textbf{(1 punto)} Reto: resuelve y comprueba.
        \[
          \dfrac{2(x - 1)}{3} - \dfrac{3(x + 2)}{4} = \dfrac{5x - 1}{6} - 2
        \]

\end{enumerate}

\vspace{10pt}
\begin{cuadroinfo}
  \small
  \textbf{Recuerda:} respeta la jerarquía de las operaciones, simplifica toda fracción hasta su forma irreductible y, en las ecuaciones, sustituye el valor obtenido para verificar la igualdad. \textquestiondown Mucho éxito!
\end{cuadroinfo}

\end{document}
